Table~\ref{tab:summary} summarises the 60 runs; Table~\ref{tab:matrix} shows every cell.

\begin{table}[t]
\centering\small
\caption{Extraction benchmark with the original prompt: number of documents (out of 6) per model and input format that were \emph{valid} (accepted by the deterministic checks), \emph{correct} (equivalent to gold), \emph{ran} (the generated solver compiled and ran) and \emph{exact} (solver output identical to the gold solver's).}
\label{tab:summary}
\begin{tabular}{@{}lcccccccc@{}}\toprule
& \multicolumn{4}{c}{HTML (MathML)} & \multicolumn{4}{c}{PDF text} \\ \cmidrule(lr){2-5}\cmidrule(lr){6-9}
LLM & valid & correct & ran & exact & valid & correct & ran & exact \\ \midrule
Llama 3 8B & 3 & 2 & 2 & 2 & 4 & 0 & 0 & 0 \\
Qwen3-Coder 30B & 2 & 1 & 1 & 1 & 3 & 2 & 2 & 1 \\
Devstral Small 2 24B & 3 & 1 & 2 & 1 & 3 & 2 & 2 & 2 \\
Qwen3.6 27B (thinking) & 6 & 4 & 3 & 2 & 5 & 3 & 3 & 2 \\
Laguna XS 2.1 33B (thinking) & 2 & 2 & 2 & 2 & 2 & 2 & 2 & 2 \\
\midrule
Total (of 30 each) & 16 & 10 & 10 & 8 & 17 & 9 & 9 & 7 \\
\bottomrule\end{tabular}
\end{table}

\begin{table}[t]
\centering\small
\caption{Outcome of every run with the original prompt (H = HTML/MathML, P = PDF text). \cmark: correct (equivalent to gold); $\sim$: accepted by the validators but not equivalent to the paper; \xmark: rejected at the trust boundary (no valid model/instance).}
\label{tab:matrix}
\begin{tabular}{@{}lcccccccccc@{}}\toprule
Document & \multicolumn{2}{c}{Llama 3 8B} & \multicolumn{2}{c}{Qwen3-Coder 30B} & \multicolumn{2}{c}{Devstral Small 2 24B} & \multicolumn{2}{c}{Qwen3.6 27B} & \multicolumn{2}{c}{Laguna XS 2.1 33B} \\
 & H & P & H & P & H & P & H & P & H & P \\ \midrule
Heat (NAFEMS) & \cmark & \xmark & \xmark & \xmark & \cmark & \cmark & \cmark & \cmark & \cmark & \cmark \\
Poisson & \cmark & $\sim$ & \cmark & \cmark & $\sim$ & \cmark & \cmark & \cmark & \cmark & \cmark \\
Transient heat & \xmark & $\sim$ & $\sim$ & $\sim$ & $\sim$ & $\sim$ & $\sim$ & $\sim$ & \xmark & \xmark \\
Elasticity & \xmark & $\sim$ & \xmark & \xmark & \xmark & \xmark & \cmark & \xmark & \xmark & \xmark \\
Stokes & $\sim$ & $\sim$ & \xmark & \cmark & \xmark & \xmark & \cmark & \cmark & \xmark & \xmark \\
Allen--Cahn & \xmark & \xmark & \xmark & \xmark & \xmark & \xmark & $\sim$ & $\sim$ & \xmark & \xmark \\
\bottomrule\end{tabular}
\end{table}

\paragraph{Overall.} Of 60 runs, 33 (55\%) produced a model and an instance that passed the deterministic checks and 19 (32\%) were equivalent to the gold specification. The two input formats are indistinguishable at this sample size (valid 16 vs.\ 17, correct 10 vs.\ 9). The 27B reasoning model (Qwen3.6) was the most reliable: valid output in 11 of 12 runs and correct output in 7; the small non-reasoning models were valid in 5 to 7 of 12. No model solved the transient-heat or the Allen--Cahn document, and the vector-valued documents were solved only by the 27B model and, once, by Qwen3-Coder.

\paragraph{The repair loop matters.} Of the 33 valid extractions, 16 (48\%) required at least one repair round, and 7 of the 19 correct ones (37\%) would not have existed without it. The repair messages were of three kinds: the JSON failed the schema (a missing \code{name}, \code{base} or \code{space} key); the equation sides had different tensor ranks; or a hallucinated dimension made the two sides of an equation inconsistent.

\paragraph{By problem class.} The steady scalar problems are largely solved: of the 20 runs on the heat and Poisson documents, 15 were correct (7 and 8). The transient-heat and Allen--Cahn documents were correct in none of the 20 runs, the elasticity document in one, the Stokes document in three. This is the pattern expected if the obstacle is the expressiveness of the prompt contract rather than mathematical understanding, which Section~\ref{sec:prompt} tests.

\subsubsection{Failure analysis}
Table~\ref{tab:failures} classifies the 41 runs that were not correct.
\begin{table}[t]
\centering\small
\caption{What happened in the 41 runs (original prompt) that were not correct.}
\label{tab:failures}
\begin{tabular}{@{}lr@{}}\toprule
Outcome of a non-correct run & runs \\ \midrule
\multicolumn{2}{@{}l}{\emph{Rejected at the trust boundary}} \\
\quad Schema/codec error (missing or mistyped key) & 16 \\
\quad Rank mismatch (validator) & 2 \\
\quad Dimension mismatch (validator) & 3 \\
\quad Output truncated (runaway reasoning) & 5 \\
\quad Parameter value missing (validator) & 1 \\
\multicolumn{2}{@{}l}{\emph{Accepted by the validators, but not equivalent to the paper}} \\
\quad Equation not equivalent to gold & 7 \\
\quad Initial condition wrong or missing & 2 \\
\quad Boundary data wrong (value, type or side) & 5 \\
\bottomrule\end{tabular}
\end{table}

\paragraph{Rejected at the trust boundary (27 runs).} The most frequent reason (16) is the \emph{schema}: the model's JSON misses a key that the codec requires (for a constitutive relation, a \code{Pow} without \code{base}, a \code{Field} without \code{space}) or contains members that do not exist in the schema (invented \code{provenance} objects). Five runs, all by the 33B reasoning model, ended in \emph{runaway reasoning} (about $62\,000$ output tokens and 20 minutes each, then truncation). Three runs by the smallest model were \emph{dimension-inconsistent} (the dimensions the model itself attached to parameters made the two sides of an equation disagree) and two were \emph{rank-inconsistent} (for example a scalar and a vector term in one sum of the Stokes momentum equation). One run lacked a parameter value ($T_\infty$ in the PDF text of the heat page, where the subscript is split from the symbol). These are the errors a typed IR detects mechanically, and in each case the diagnostic names what is wrong.

\paragraph{Accepted by the validators but not equivalent to the paper (14 runs).} This is the dangerous class. Seven runs have a wrong equation (for example a flipped source sign, a missing time-derivative term, an equation that refers to the time as an unknown field, or the wrong momentum operator), five have wrong boundary data (a ramp written with a wrong time symbol; a prescribed temperature turned into a Robin condition), and two lost the initial condition (the Allen--Cahn document has a spatially varying initial field that the original prompt gives the model no node to express). Eight of the transient-heat extractions write the ramp as \code{beta*t} with \code{t} as a declared \emph{field} or \emph{parameter} instead of the time coordinate (Appendix~\ref{app:supp}); the original prompt does not list a time coordinate among its node tags.

\subsubsection{End-to-end consequences}
\label{sec:e2e}
The 19 correct extractions entered the pipeline (generation, build, run; verification disabled) with the discretization data of the gold instance. Eighteen compiled and ran, and 15 of these produced output \emph{identical} to the gold solver's (relative nodal $\ell^2$ difference below $10^{-6}$). Three runs on the Stokes document were correct only \emph{modulo the incompressibility constraint}: the models wrote $-\nabla\cdot(\nu\nabla u)+\nabla p$ instead of $-\nabla\cdot(2\nu\varepsilon(u))+\nabla p$. The solvers ran, and their solutions differ from the gold one by 16\% (velocity) and 9.6\% (pressure) in relative nodal $\ell^2$ norm: the two forms agree in the continuum for $\nabla\cdot u=0$ but are different discrete problems (a different treatment of the lid corners is consistent with the velocity difference as well). One further run, correct by the scorer, was refused by the generator: the Navier--Cauchy form of elasticity (Appendix~\ref{app:supp}).

What stopped the 14 ``valid but wrong'' runs? Thirteen did not produce a solution. In eight the generator refused the model (``field \code{t} uses space L2'', ``Dirichlet value must be a Constant or depend on space/time'', ``a time derivative but no initial condition''); in three the pipeline's \code{validate} stage rejected it (one boundary condition was flagged by the anti-dissipative flux rule; two instances failed because the injected pressure pinning referred to a field the model had named differently, an artefact of our injection of the gold discretization data); one produced code that failed to compile; and one made the generator crash. Only \emph{one} of the 14 ran to completion with a wrong problem: the Poisson source of flipped sign (Appendix~\ref{app:supp}). Two points should be made plainly. First, the late layers worked here largely because the backend's scope is narrow, not because they understand the mathematics: the generator's refusals are a by-product of what it cannot represent, so a wider backend would let more of these through, and the compile failure and the crash show that the generator's robustness against malformed models is itself a gap. Second, errors such as the flipped source sign are, by construction, invisible to every internal consistency check; they can only be detected by human review or by comparison with an independent reference, which is why the framework stores what was extracted together with the source notes.

\subsubsection{Cost}
Table~\ref{tab:cost} gives the average cost per run (two to three model calls including repair rounds). The non-reasoning models answer in well under a minute per call after loading: the median run takes 46~s (Llama~3), 40~s (Qwen3-Coder) or 3.3 minutes (Devstral, which is slower on this server). The two reasoning models spend 85\% and 96\% of their output on reasoning, which makes a typical Qwen3.6 run take 6 minutes and a bad Laguna run 20 minutes. The prompt of a first attempt is about 3\,000 tokens for the model call and 2\,000 for the instance call; the totals in the table include the repeated context of repair rounds. On the heat page, a correct extraction cost about $1.2\cdot10^3$ output tokens with the non-reasoning models and about $9\cdot10^3$ with Qwen3.6. No comparison with direct code generation is made, since that baseline was not run.
\begin{table}[t]
\centering\small
\caption{Mean cost per run (original prompt, 12 runs per model): LLM calls, prompt and output tokens (summed over the calls of a run, including repair rounds), fraction of the output that is reasoning, and wall time per run on the local server.}
\label{tab:cost}
\begin{tabular}{@{}lrrrrrr@{}}\toprule
LLM & calls & prompt tok. & output tok. & reasoning share & median s/run & max s/run \\ \midrule
Llama 3 8B & 3.1 & 7\,951 & 1\,913 & 0\% & 46 & 62 \\
Qwen3-Coder 30B & 2.7 & 8\,361 & 2\,573 & 0\% & 40 & 73 \\
Devstral Small 2 24B & 2.9 & 9\,288 & 2\,845 & 0\% & 198 & 279 \\
Qwen3.6 27B (thinking) & 2.9 & 7\,207 & 13\,260 & 85\% & 343 & 859 \\
Laguna XS 2.1 33B (thinking) & 2.0 & 5\,946 & 30\,718 & 96\% & 124 & 1235 \\
\bottomrule\end{tabular}
\end{table}

\subsection{Experiment: how much does the prompt contract matter?}
\label{sec:prompt}
The failure analysis suggests that many errors arise because the original extraction prompt does not tell the model how to express things the IR \emph{can} express: it lists node tags without the time and space coordinates, vector construction, initial conditions or the form of a constitutive relation, and its only format example is a scalar Poisson model. To test this hypothesis we extended the system prompt by a short block (Appendix~\ref{app:supp}; about 350 tokens) that documents exactly these node types and rules, and re-ran the six documents in HTML for the four models that finished in the original experiment (the runaway-reasoning model was not repeated). Nothing else changed: same documents, same temperature, same scorer.

\begin{table}[t]
\centering\small
\caption{Effect of the extended prompt contract (HTML documents, six per model): documents valid, correct and ran (solver compiled and ran) with the original and with the extended prompt.}
\label{tab:v2}
\begin{tabular}{@{}lcccccc@{}}\toprule
& \multicolumn{3}{c}{original prompt} & \multicolumn{3}{c}{extended prompt} \\ \cmidrule(lr){2-4}\cmidrule(lr){5-7}
LLM & valid & correct & ran & valid & correct & ran \\ \midrule
Llama 3 8B & 3 & 2 & 2 & 4/6 & 0/6 & 0/6 \\
Qwen3-Coder 30B & 2 & 1 & 1 & 5/6 & 4/6 & 3/6 \\
Devstral Small 2 24B & 3 & 1 & 2 & 4/6 & 4/6 & 3/6 \\
Qwen3.6 27B & 6 & 4 & 3 & 6/6 & 6/6 & 4/6 \\
\bottomrule\end{tabular}
\end{table}

\begin{table}[t]
\centering\small
\caption{Per-document outcome with the extended prompt (\cmark\ correct, $\sim$ valid but not equivalent, \xmark\ rejected at the trust boundary).}
\label{tab:v2matrix}
\begin{tabular}{@{}lcccccc@{}}\toprule
LLM & Heat (NAFEMS) & Poisson & Transient heat & Elasticity & Stokes & Allen--Cahn \\ \midrule
Llama 3 8B & $\sim$ & $\sim$ & $\sim$ & \xmark & $\sim$ & \xmark \\
Qwen3-Coder 30B & \cmark & \cmark & \cmark & \cmark & \xmark & $\sim$ \\
Devstral Small 2 24B & \cmark & \cmark & \cmark & \xmark & \cmark & \xmark \\
Qwen3.6 27B & \cmark & \cmark & \cmark & \cmark & \cmark & \cmark \\
\bottomrule\end{tabular}
\end{table}

Across the four models, the number of documents (out of 24) that were valid rose from 14 to 19, the number of correct ones from 8 to 14 (from 33\% to 58\%), and the number whose solver compiled and ran from 8 to 10 (Table~\ref{tab:v2}). The effect is concentrated where the original prompt had a gap (Table~\ref{tab:v2matrix}). The transient-heat document, solved by none of the models before, was now correct for three of four, and the three corresponding solvers ran and reproduced the gold output exactly: the models used the time coordinate. For Allen--Cahn the 27B model now extracted the initial field and the relation, and its solver output equals the gold output to $4\cdot10^{-17}$. Qwen3-Coder and the 27B model now produced a correct elasticity model. The effect is not uniform: the smallest model got \emph{worse} (from 2 to 0 correct): with a prompt about twice as long it produced equations that fail the scorer and, in two runs, programs whose generated C++ did not compile, so a longer contract is not automatically better for an 8B model.

\paragraph{The contract has to match the backend.} Only 10 of the 14 correct extractions ran. The reason is instructive: our own extension told the models to write a zero vector as a \code{VectorExpr} of zeros, which the generator rejects (it accepts a \code{Constant} $0$ for the zero vector), and several models declared the pressure in $L^2$, which the generator does not support. Two of the four correct elasticity and Stokes extractions were therefore refused by the \emph{backend}, i.e., the prompt promised more than the generator delivered. We corrected both points (a ``v3'' prompt: zero vector as \code{Constant} $0$; all fields declared \code{H1}, the spaces being chosen by the solver) and re-ran the two vector-valued documents with three models (Table~\ref{tab:v3}). Both Stokes extractions that were correct now ran (they are correct only modulo the incompressibility constraint, with a 16\% velocity and a 33\% pressure difference to the gold discretization, the same kind of difference as in Section~\ref{sec:e2e}); the 27B model's elasticity extraction was correct but the generated program did not compile (the generator emitted invalid C++ for that equivalent formulation), and Qwen3-Coder's Stokes extraction was accepted by every check, ran, and solved a different problem (a pressure that differs from the gold one by more than 100\%), a second silently wrong run in the whole study. The lesson is that the prompt contract, the IR schema and the generator's actual scope should be derived from one specification; today they are three hand-maintained texts.

\begin{table}[t]
\centering\small
\caption{Second revision of the prompt (v3; elasticity and Stokes only, HTML). \cmark\ correct, $\sim$ valid but not equivalent, \xmark\ rejected; ``ran'' and ``exact'' count documents whose solver compiled and ran, and whose output equals the gold output.}
\label{tab:v3}
\begin{tabular}{@{}lcccc@{}}\toprule
LLM & Elasticity & Stokes & ran (of 2) & exact (of 2) \\ \midrule
Qwen3-Coder 30B & $\sim$ & $\sim$ & 1 & 0 \\
Devstral Small 2 24B & \xmark & \cmark & 1 & 0 \\
Qwen3.6 27B (thinking) & \cmark & \cmark & 1 & 0 \\
\bottomrule\end{tabular}
\end{table}

\paragraph{Caveat.} The extended prompts were written \emph{after} the original failures were analysed and evaluated on the same documents. The improvement shows that the prompt contract is a first-order lever, and it is probably an upper bound for what a generic revision would achieve on new documents. A fair evaluation requires a held-out corpus (Section~\ref{sec:outlook}).



\subsection{A toy direct-generation baseline}
\label{sec:baseline}
The research plan names ``direct LLM $\to$ \dealii{}'' as the reference architecture (baseline A). We did not run the full ablation, but we ran its smallest meaningful instance, on the heat page of Section~\ref{sec:heat}, to calibrate the claim that direct generation is hard for \dealii{}. Three of the locally served models were asked to write a complete \code{main.cc} from the same page text (continuous Lagrange elements on a structured mesh, print the value at the reference point, free choice of the mesh), with up to four attempts; after a failed attempt the model received the first compiler errors (or the run failure), the feedback protocol of PDEAgent-Bench \citep{pdeagentbench2026}. A run succeeds if the printed value at $P$ is within $10^{-2}$ of $18.25$. The build used the same \dealii{} installation as all other experiments.

\begin{table}[t]
\centering\small
\caption{Direct generation of a \dealii{} program for the heat page (one document, temperature $0$, at most four attempts with compiler feedback). Tokens are summed over all calls.}
\label{tab:baseline}
\begin{tabular}{@{}lccrrr@{}}\toprule
LLM & attempts & outcome & probe error & tokens (in+out) & wall time \\ \midrule
Devstral Small 2 24B & 4 & no compile & -- & 22\,327+8\,995 & 697\,s \\
Qwen3-Coder 30B & 4 & no compile & -- & 16\,006+5\,036 & 183\,s \\
Qwen3.6 27B (thinking) & 3 & no compile (truncated) & -- & 11\,389+27\,168 & 884\,s \\
\bottomrule\end{tabular}
\end{table}

None of the three models produced a program that compiled (Table~\ref{tab:baseline}). The failures are the typical API errors of a large C++ library: constructing a \code{SparseMatrix} directly from a \code{DynamicSparsityPattern}, calling a non-const member of a const \code{FEValues}, a wrong \code{Quadrature} constructor, and include files that do not exist in this version of the library (\code{constraints.h}, \code{matrix\_tools.h}); the repair rounds repeated the same class of error in three of the four attempts for two of the models, and the reasoning model exhausted its budget of $12\,000$ output tokens on the third attempt. For the same page and models, the \femagent{} extraction needed one model call and one instance call, about $6\cdot10^3$ tokens for Devstral and $1.1\cdot10^4$ for Qwen3.6, and produced a valid, correct specification in one attempt with both models; Qwen3-Coder failed this page with the original prompt (a schema error in three attempts) and passed it with the extended one.

This comparison is deliberately small and favours no one strongly: it is one document and one sample per model, our direct-generation prompt is minimal (no library API guide, no template, which PDEAgent-Bench shows to help), and stronger or larger models would very likely do better; in that benchmark the best single-shot pass rate on \dealii{} is $32.3\%$ even for frontier models. What it does show is the \emph{kind} of failure the architecture avoids: the unit of failure is a missing JSON key, which is named in one line and repaired cheaply, not a compiler error in a library API that the model does not know.

